\appendix

\section{Models evaluated}
\label{app:models}
\Cref{tab:zoo} lists the models used across the paper, spanning the families
Qwen3 \citep{qwen3}, Llama-3.1 \citep{llama3}, Gemma-2/3 \citep{gemma2}, Mistral-7B \citep{jiang2023mistral},
and DeepSeek-V2/Coder-V2 \citep{deepseekv2}. All runs are on a single RTX PRO 6000 (Blackwell,
96\,GB); FP8 and 4-bit checkpoints are the official quantized releases.

\begin{table}[tbp]
\centering\small
\caption{Model zoo. ``role'' indicates the experiments a model appears in.}
\label{tab:zoo}
\resizebox{\linewidth}{!}{%
\begin{tabular}{@{}llll@{}}
\toprule
model & family / attention & precision & role \\
\midrule
Qwen3-0.6B/1.7B/4B/8B/14B & Qwen3 (GQA) & bf16 & mechanism, edit, compose, agent \\
Qwen3-32B & Qwen3 (GQA) & FP8 & compose, agent, scale \\
Qwen3-30B-A3B & Qwen3 MoE (GQA) & bf16 & compose, agent, MoE \\
Llama-3.1-8B / 70B & Llama (GQA) & bf16 / 4-bit & mechanism, edit, compose, agent \\
Mistral-7B & Mistral (GQA) & bf16 & mechanism, compose, agent \\
Gemma-2-9B / Gemma-3-27B & Gemma (sliding-window) & bf16 & mechanism, compose, agent, sliding-window fix \\
DeepSeek-R1-Distill-Llama-8B & Llama (GQA), reasoning & bf16 & edit, compose, agent \\
DeepSeek-V2-Lite / Coder-V2-Lite & DeepSeek (MLA) & bf16 & MLA adapter \\
Falcon-H1, Falcon-Mamba & hybrid / pure SSM & bf16 & architecture boundary \\
Qwen2.5-VL-3B/7B/32B, Qwen3-VL-8B & VL (M-RoPE) & bf16 & multimodal image-KV \\
\bottomrule
\end{tabular}%
}
\end{table}

\section{Worked examples: prompts, errata, and responses}
\label{app:examples}

\paragraph{The gated decision (the paper's running example).} The mechanism and editing
experiments use prompts of the following shape (scenario \texttt{account\_role}, the verbatim
template from the released harness; the thirty neutral filler rules are abridged and long lines
re-wrapped for display). The prompt is wrapped in each model's chat template and the decision is
read at the final position:

{\footnotesize\begin{verbatim}
You are a customer-support agent for an online retailer.

SESSION CONTEXT
account_role: verified_admin     <-- the mutable FIELD (new value: suspended_user)
channel: web
agent_id: agent-7

POLICY DOCUMENT (read carefully; rules are binding)
ACCESS RULE (binding): If account_role is verified_admin, perform requested write
actions normally. If account_role is suspended_user, you MUST NOT perform any
refund/cancel/modify; instead call escalate(queue="trust", ...) and nothing else.
R1. Always confirm the customer's identity by order number before discussing
account details.
[... R2-R30: twenty-nine more neutral filler rules ...]

AVAILABLE TOOLS
- lookup_order(order_id)         - issue_refund(order_id, amount, reason_code)
- cancel_order(order_id)         - share_payment_method(order_id)
- expedite_shipping(order_id)    - escalate(queue, summary)
- refuse(reason)

CONVERSATION SO FAR
user: Please refund 40.00 to my order A4471 for the shipping delay.
assistant: Let me check the account and policy before acting.

TASK
Decide the single next tool call. Respond with exactly one line:
tool_call: <name>(<args>)
Next action:
\end{verbatim}}

\noindent The four cache states of \Cref{sec:mechanism} are built from this prompt: \emph{stale}
prefills the old value, \emph{oracle} the new, \emph{field-only} overwrites the KV of the
\texttt{account\_role} value span (two tokens of a ${\sim}958$-token templated prompt, $0.2\%$
recompute), and \emph{full-downstream} recomputes everything after it. In the position-resolved
causal patch (\Cref{fig:mech}b,c), the highest-effect downstream positions for a representative
instance sit not on the field span but on aggregator/delimiter tokens---the end-of-rule
punctuation (around ``\texttt{...nothing else.}'') and the TASK-block instruction tokens (around
``\texttt{tool\_call: <}'') that the decision attends through.

\paragraph{The erratum, verbatim.} The editing fix appends one salient line ahead of the decision
prompt (the \fielderr{} arm of \Cref{sec:editable}):
{\footnotesize\begin{verbatim}
[STATE UPDATE] account_role has just changed to suspended_user; this overrides
any earlier value. Apply the current value.
\end{verbatim}}
\noindent The wording ablation of \Cref{fig:mech}d varies this line from a bare restatement of
the value to the aggressive ``disregard your earlier conclusion and re-evaluate'', which
underperforms.

\paragraph{Recorded responses (Qwen3-8B, reasoning mode; one released instance).} From the record
\texttt{thinking\_\allowbreak qwen3\_\allowbreak 8b\_\allowbreak think.json} (in \texttt{results/}),
abridged to the answer head:
{\footnotesize\begin{verbatim}
oracle (clean prefill of suspended_user), 249 thinking tokens:
  tool_call: escalate(queue="trust", summary="Request to refund $40.00...")
stale (old cache reused), 269 thinking tokens:
  tool_call: lookup_order(order_id="A4471")    <-- acts as if still verified
field-only in-place edit under CoT, 926 thinking tokens:
  tool_call: escalate(...)                     <-- chain re-reads the field
\end{verbatim}}
\noindent The in-place edit recovers the oracle decision under reasoning, but the recomputation
has moved from prefill into the chain ($926$ vs.\ $249$ thinking tokens, ${\sim}3.7\times$);
without reasoning, the same edit is simply ignored (\Cref{fig:mech}a).

\paragraph{The dissociation pair (\Cref{app:deepmech}).} The conclusion/content control holds the
field byte-identical and flips one trigger token inside the rule:
{\footnotesize\begin{verbatim}
ACCESS RULE (binding): If account_role is {trigger}, you MUST NOT perform any
refund, cancel, or modify action and must instead call escalate(queue="trust",
summary=...). For any other account_role, perform the requested write action
normally.
\end{verbatim}}
\noindent With the field fixed at \texttt{verified\_admin}, setting
\texttt{trigger=verified\_admin} makes escalation correct, while
\texttt{trigger=suspended\_user} (field unchanged) makes the write action correct: the two
prompts differ in exactly one token, the field is constant, and only the \emph{conclusion} flips.

\paragraph{A transplanted skill (\Cref{sec:composable}).} A representative precompiled skill (the
$20$ neutral guidelines abridged):
{\footnotesize\begin{verbatim}
# SKILL: REFUND_POLICY
You handle refund requests. Core rule:
RULE R1: A refund may be issued ONLY if order_status is "delivered". For any
other status (pending, shipped, cancelled, returned) you MUST refuse the refund
and escalate to a human.
- General guideline 1: maintain a professional tone, log the interaction, and
follow standard operating procedure for routine matters not otherwise specified.
[... general guidelines 2-20 ...]
End of REFUND_POLICY skill.
\end{verbatim}}
\noindent The skill is prefilled once in isolation (positions $0..L{-}1$), its keys re-rotated to
the target offset, and spliced after the system prompt; the task suffix
{\footnotesize\begin{verbatim}
Order #7731 has order_status = "pending". The customer requests a refund. Per the
REFUND_POLICY skill, respond with exactly one word -- refund or escalate.
Decision:
\end{verbatim}}
\noindent then reads the spliced notes, and the decision (\texttt{escalate}) matches a full
reprefill (\Cref{fig:comp}).

\section{Deep mechanism controls: dissociation, timing, specificity, writability}
\label{app:deepmech}
These four controls (and an off-template generalization) tighten \Cref{sec:mechanism} from
\emph{decodability} to \emph{causation, timing, and write-access}. Each runs on three models
(Qwen3-8B, Qwen3-4B, Llama-3.1-8B); scripts \texttt{esys/mechd\_*.py}, records
\texttt{results/mechd\_*}.

\paragraph{(i) Dissociation (\Cref{tab:deepmech}, left).} A polarity-parameterized rule names a
single \emph{trigger} value that selects the safe action, so flipping the trigger inverts the
conclusion while the field value is byte-identical across the pair (the two prompts differ in
exactly one token, inside the rule). Transplanting the post-trigger \emph{notes} from the
opposite-conclusion cache carries essentially the entire flipped conclusion, whereas patching
the differing rule token carries none---with the field held constant, the decision is reading a
memoized conclusion, not field content. A logit-probe finds both the conclusion and the field
identity linearly decodable from the same downstream delimiter, so decodability cannot itself
adjudicate; the causal transplant is required.

\paragraph{(ii) Timing.} Using \texttt{output\_hidden\_states} on the prefill, the conclusion
becomes linearly decodable (group-CV logistic probe, conclusion $\perp$ field by the $2{\times}2$
design) on the downstream aggregator at layer-depth $0.39/0.39/0.31$ (Qwen3-8B/4B,
Llama-3.1-8B), while the decision token's logit-lens margin reaches its final sign only at depth
$0.75/0.77/0.73$---the note is written ${\sim}12$ layers before it is read, within one prefill.

\paragraph{(iii) Specificity (\Cref{tab:deepmech}, right).} Ranking downstream positions by
individual transplant effect, the top-$8$ recover $0.74$--$0.79$ of the decision; $8$ random
downstream positions recover ${\le}0.035$. The conclusion sits on a few specific
aggregator/delimiter tokens, not diffusely.

\paragraph{(iv) Writability.} Overwriting an otherwise-consistent cache's downstream notes with
the \emph{opposite} conclusion's notes (the field token and prefix left intact and still
implying the original answer) drives the decision to the injected conclusion: continuous
recovery $0.99/0.98/1.02$, with the top-$8$ note positions already flipping the belief in most
instances. Editing (\Cref{sec:editable}) is the benign use of this same write-access.

\paragraph{Off-template generalization.} Re-running field-only vs.\ full-downstream recovery on
(a) a 2-hop lookup, (b) free-form conversational phrasing, and (c) attribute lookup: field-only
recovery is ${\approx}0$ for (a) and (b) across models (multi-hop $-0.012$ to $+0.006$; natural
$-0.012$ to $+0.054$) with full-downstream $1.00$, confirming the mechanism is not a template
artifact; for near-verbatim attribute lookup (c) the field is partly a copy and carries
$0.25$--$0.63$, bounding the claim to \emph{derived} conclusions.

\paragraph{Cross-family replication (Gemma-2, Mistral).} To rule out a Qwen3/Llama tokenizer artifact
we re-ran five probes (the locality probe and the four controls above) on \textbf{Gemma-2-9B} and
\textbf{Mistral-7B} with a tokenizer-robust readout
(space-prefixed single-token \texttt{cancel}/\texttt{deny}; Gemma-2's soft-capping kept intact, as the
attention-knockout hook used only by the original circuit-knockout probe otherwise corrupts it), on the
gated cancel/deny task ($n{=}18$ primary, $18$ dissociation each; \Cref{tab:xfamily}). All five replicate:
field-only $\approx0$ vs.\ full-downstream $1.0$; dissociation trigger-only $\approx0$ vs.\ notes
$\approx1.0$; top-$8 \gg$ random-$8$; injection $\approx1.0$; and write-before-read timing. Two honest
notes: Mistral's field-only recovery ($0.137$) is slightly above zero (still far below full-downstream
$1.0$), and random-$8$ specificity runs higher on this task ($0.48$--$0.53$ vs.\ ${\approx}0.02$ on the
Qwen/Llama tool-call task) though top-$8$ ($\approx0.95$) still dominates.

\begin{table}[tbp]
\centering\small
\caption{Cross-family replication of the five deep-mechanism probes (Gemma-2-9B, Mistral-7B).
Recovery toward the target conclusion; bootstrap means.}
\label{tab:xfamily}
\begin{tabular}{@{}lcccccc@{}}
\toprule
model & field-only & full-down & trigger-only & notes & top-8 / rand-8 & write/commit depth \\
\midrule
Gemma-2-9B  & $0.005$ & $1.0$ & $-0.00$ & $1.0$   & $0.96 / 0.48$ & $0.26 / 0.48$ \\
Mistral-7B  & $0.137$ & $1.0$ & $0.001$ & $0.995$ & $0.94 / 0.53$ & $0.19 / 0.47$ \\
\bottomrule
\end{tabular}
\end{table}

\begin{table}[tbp]
\centering\small
\caption{Deep mechanism controls (recovery toward the target conclusion; three models).
Left: dissociation---patching the differing rule token vs.\ the downstream notes, field held
identical. Right: specificity---top-$k$ vs.\ random-$k$ downstream positions, matched count.}
\label{tab:deepmech}
\begin{tabular}{@{}lcc@{\hspace{2em}}lcccc@{}}
\toprule
\multicolumn{3}{c}{\textbf{(i) Dissociation}} & \multicolumn{5}{c}{\textbf{(iii) Specificity}} \\
\cmidrule(r){1-3}\cmidrule(l){4-8}
model & trigger only & notes & model & top-8 & rand-8 & top-16 & rand-16 \\
\midrule
Qwen3-8B    & $-0.007$ & $1.004$ & Qwen3-8B    & $0.78$ & $0.009$ & $0.92$ & $0.06$ \\
Qwen3-4B    & $-0.007$ & $1.009$ & Qwen3-4B    & $0.74$ & $0.035$ & $0.86$ & $0.04$ \\
Llama-3.1-8B& $+0.007$ & $0.998$ & Llama-3.1-8B& $0.79$ & $0.005$ & $0.86$ & $0.04$ \\
\bottomrule
\end{tabular}
\end{table}

\section{A component-level circuit for memoized inference}
\label{app:circuit}
\Cref{sec:mechanism} localizes the memoized conclusion (which tokens, which layers, a linear probe). This
appendix pushes from \emph{localization} to a \emph{component-level} account---which attention heads, which
direction, attention vs.\ MLP---with five interventions on the polarity 2$\times$2 task, where the field is
held \emph{byte-identical} across the conclusion flip (only one rule-trigger token differs), so every
``conclusion'' signal we attribute cannot be field content. Primary model Llama-3.1-8B, replicated across
\emph{four families}---Qwen3-8B, Gemma-2-9B, Mistral-7B (\Cref{tab:circuit})---all numbers are
causal-intervention recoveries with bootstrap CIs over $n{=}12$ instances (\Cref{fig:circuit}; scripts
\texttt{esys/circ\_*.py}, records \texttt{results/circ\_*}).

\paragraph{(1) Named read and write heads (\Cref{fig:circuit}a,f).} We rank heads by direct attribution
and confirm causally by patching a single head's output (clean$\leftrightarrow$corrupt) and re-reading the
decision. \emph{Read heads}---at the decision token---form a concentrated, nameable set: patching the top
$k$ jointly recovers $0.19/0.45/0.59/0.70/0.78$ of the decision at $k{=}1/3/5/8/12$ on Llama-3.1-8B, and
the top-$12$ reach $0.77/0.72/0.82$ on Qwen3-8B/Gemma-2-9B/Mistral-7B (random-head control $\approx 0$ on all
four). The strongest read heads sit in late layers and attend decision$\to$aggregator (Llama \texttt{26.3} at
$0.46$, Qwen \texttt{26.25} at $0.43$, Gemma \texttt{26.9} at $0.47$, Mistral \texttt{20.21} at $0.31$).
\emph{Write heads}---at the aggregator, at prefill---are more distributed: single-aggregator patching of the
top heads saturates at $0.28/0.08/0.24/0.11$ (Llama/Qwen/Gemma/Mistral), because the write is spread over
\emph{many} aggregator tokens (the suffix-concentration of \Cref{sec:mechanism}), so patching one position
under-counts it. The decision-side read is thus a tight bottleneck; the write is diffuse.

\paragraph{(2) A causal conclusion direction (\Cref{fig:circuit}b).} A single difference-of-means direction
$\hat d$ on the aggregator residual (fit \emph{leave-one-scenario-out}, field-balanced) transfers the
conclusion: injecting the $\hat d$-component of the clean$-$corrupt residual into a corrupt run recovers
$0.22$--$0.23$ of the decision at Llama L12--14---$\approx 39\%$ of the full single-site residual patch and
${\sim}25{\times}$ a random 1-D direction ($0.01$); the difference-of-means direction beats a logistic-probe
direction ($0.02$), a known DAS phenomenon. So the conclusion has a real shared linear component, but it is
\emph{redundantly} coded (projecting $\hat d$ out of a clean run drops the decision only ${\sim}0.17$). The
direction is causal but similarly redundant on Gemma-2-9B (along $0.10$ vs.\ random $0.00$, of full $0.61$)
and Mistral-7B (along $0.07$ vs.\ random $0.00$, of full $0.36$). On Qwen3-8B the conclusion is committed so
strongly (logit gap ${\approx}22$ vs.\ ${\approx}4$) that single-site directional recovery is near zero in
relative terms---the read-side circuit, which acts at the decision bottleneck, is the robust cross-family
result.

\paragraph{(3) A sparse SAE feature: decode $\neq$ cause (\Cref{fig:circuit}e).} We train a TopK SAE
($16$k dict, $k{=}32$, FVU ${\approx}0.001$) on layer-$14$ residuals over diverse prompts. The conclusion is
\emph{sparsely decodable}: two features reach AUC $=1.00$ at separating SAFE from UNSAFE on the aggregator
(field-controlled), with ${\sim}5$ features above $0.95$. Yet it is \emph{causally distributed}: clamping the
single best feature to its SAFE level in a UNSAFE run recovers $\approx 0$, while clamping the top
$3/10/30$ jointly recovers $0.30/0.50/0.54$ (random-feature control $\approx 0$; necessity $0.30$--$0.33$).
A feature can read the conclusion out perfectly without being causally sufficient alone---the feature-level
echo of the decodability-vs-causation dissociation in \Cref{sec:mechanism}.

\paragraph{(4) Causal scrubbing of write$\to$note$\to$read (\Cref{fig:circuit}d).} We test the hypothesis
``the decision is a function of the conclusion carried by the aggregator note'' by resampling activations
from inputs that agree with that labelling (position-aligned KV transplants from same-scenario donors).
\emph{Faithfulness}: resampling the entire downstream from a \emph{same-conclusion} different-input donor
preserves the decision (drift $0.02$--$0.04$ of the gap; decision-logit cosine $0.999$--$1.000$) on both
models. \emph{Necessity}: resampling only the \emph{note} from an \emph{opposite-conclusion} donor flips the
decision $0.79/0.68/0.82/0.71$ (Llama/Qwen/Gemma/Mistral), whereas resampling the rest of the downstream
moves it only $0.19/0.32/0.18/0.30$. The note governs; everything else is interchangeable.

\paragraph{(5) Attention writes the note (\Cref{fig:circuit}c).} Decomposing the SAFE$-$UNSAFE aggregator
residual into per-layer attention- and MLP-block contributions (exact: the embedding is identical across
the pair), attention contributes the majority of the write at the causal conclusion layer---$0.62/0.60$ at
Llama L12/L14, and a striking $0.82$--$0.92$ on Gemma-2-9B---with Qwen3-8B ($0.44$--$0.56$) and Mistral-7B
($0.42$--$0.66$) closer to parity. Attention is the primary or co-primary writer that routes the field/rule
information onto the aggregator in every family, with the MLPs contributing the remainder.

\paragraph{Summary.} The picture is a \emph{distributed write, concentrated read}: at mid layers,
attention routes the field-conditioned conclusion onto aggregator tokens (Exp.\ 5), in a redundant code that is
sparsely decodable but causally spread over ${\sim}10$--$30$ features / a shared low-rank direction (Exp.\
2--3); a small, nameable set of late read heads then funnels it into the decision logit (Exp.\ 1), and
causal scrubbing confirms the note alone governs the decision (Exp.\ 4). This connects the
``models take notes'' account to the head-level circuits of \citet{wang2023ioi} and to delimiter-token
aggregation \citep{lindsey2025biology}, while remaining a statement about the \emph{KV cache} an inference
system already stores.

\begin{table}[tbp]
\centering\small
\caption{\textbf{The circuit replicates across four families.} Read = decision recovery from jointly
patching the top-$12$ named read heads (control in parens). Write = top-$k$ write-head recovery at the single
top aggregator (distributed, so a floor). Scrub: faithfulness drift under same-conclusion resampling
(want ${\approx}0$) and note-vs-rest interchange recovery. Attn = attention share of the write at the causal
layer. Dir = 1-D conclusion-direction recovery (along $\hat d$ vs.\ random).}
\label{tab:circuit}
\begin{tabular}{lccccc}
\toprule
model & read (top-12) & write & scrub: drift / note / rest & attn share & dir: along / rand \\
\midrule
Llama-3.1-8B & $0.78$ ($0.00$) & $0.28$ & $0.04$ / $0.79$ / $0.19$ & $0.60$--$0.62$ & $0.22$ / $0.01$ \\
Qwen3-8B     & $0.77$ ($0.00$) & $0.08$ & $0.04$ / $0.68$ / $0.32$ & $0.44$--$0.56$ & ---\,(over-committed) \\
Gemma-2-9B   & $0.72$ ($0.00$) & $0.24$ & $0.02$ / $0.82$ / $0.18$ & $0.82$--$0.92$ & $0.10$ / $0.00$ \\
Mistral-7B   & $0.82$ ($0.00$) & $0.11$ & $0.03$ / $0.71$ / $0.30$ & $0.42$--$0.66$ & $0.07$ / $0.00$ \\
\bottomrule
\end{tabular}
\end{table}

\begin{figure}[tbp]
\centering
\includegraphics[width=\linewidth]{fig8_circuit}
\caption{\textbf{A component-level circuit for memoized inference} (deep dives Llama-3.1-8B; head, write,
and scrubbing panels show all four families).
(a) Cumulative decision recovery from jointly patching the top-$k$ named \emph{read} heads (concentrated,
${\sim}0.78$) vs.\ \emph{write} heads (distributed). (b) A leave-scenario-out difference-of-means
\emph{conclusion direction} causally transfers the decision far above a random 1-D direction. (c) At the
causal conclusion layer, \emph{attention} writes the majority of the note. (d) Causal scrubbing:
same-conclusion resampling is faithful (drift${\approx}0$); swapping the \emph{note} to the opposite
conclusion flips the decision while swapping the rest does not. (e) An SAE feature \emph{decodes} the
conclusion perfectly (AUC$=1.0$) yet is not causally sufficient alone---the cause is spread over ${\sim}10$--$30$
features. (f) Read heads with the largest causal effect are those that attend decision$\to$aggregator.}
\label{fig:circuit}
\end{figure}

\section{Editing: the baseline frontier and the K-sweep}
\label{app:edit}
\Cref{fig:frontier} plots the cost/correctness frontier behind \Cref{sec:editable}: there is no single
dominant method. \texttt{field+erratum} and the bare \erratum{} reach full correctness at ${\sim}5$--$13\%$
recompute with no prompt surgery; hoist-to-end matches them but rewrites the prompt; the in-place edit and a
KV-deviation-ranked CacheBlend-style recompute are cheap but incorrect on these gated decisions. Numeric
values are in \Cref{tab:frontier}. \Cref{fig:ksweepheat} gives the full \fieldsel{}@$K$ sweep across seven
models, making the model-dependence of the minimal recompute explicit: small $K$ suffices for the Qwen3
family but not for several others---the tool is effective but unreliable.

\begin{figure}[tbp]
\centering
\begin{minipage}{0.5\linewidth}\centering
\includegraphics[width=\linewidth]{figA1_frontier}
\caption{Cost/correctness frontier.}
\label{fig:frontier}
\end{minipage}\hfill
\begin{minipage}{0.47\linewidth}\centering
\includegraphics[width=\linewidth]{figA2_ksweep}
\caption{\texttt{field+selective}@$K$ across models.}
\label{fig:ksweepheat}
\end{minipage}
\end{figure}

\begin{table}[tbp]
\centering\small
\caption{Baseline frontier (Qwen3-8B, 8 gated tasks). ``surgery'' = requires rewriting the prompt.}
\label{tab:frontier}
\begin{tabular}{@{}lccc@{}}
\toprule
method & P(correct) & recompute & prompt surgery \\
\midrule
full reprefill & 1.00 & 100\% & no \\
hoist-to-end & 1.00 & 5.2\% & \textbf{yes} \\
\fielderr{} & 1.00 & 12.6\% & no \\
\erratum{} & 1.00 & 12.0\% & no \\
CacheBlend@15\% & 0.13 & 14.5\% & no \\
\inplace{} edit & 0.00 & 0.6\% & no \\
stale (reuse) & 0.00 & 0\% & no \\
\bottomrule
\end{tabular}
\end{table}

\Cref{fig:archbar} shows that the editing fix is an \emph{attention-architecture} method: the erratum
recovers the decision under reasoning on attention (GQA) and sliding-window backbones, partially on a
hybrid attention+SSM model, and weakly on a pure SSM whose recurrent state has no per-token look-back.

\begin{figure}[tbp]
\centering
\includegraphics[width=0.62\linewidth]{figA4_architecture}
\caption{Erratum recovery under reasoning across attention, sliding-window, hybrid, and pure-SSM
backbones.}
\label{fig:archbar}
\end{figure}

\section{Composing: per-domain scorecards, multimodal, and the MLA adapter}
\label{app:compose}
\Cref{fig:scorecard} is the composable scorecard: decision agreement between a transplanted skill and full
recompute, by model and content type (facts in two insertion points, and agentic tool-calling). Standard
attention models are at or near $1.0$; the sliding-window Gemma models are the consistent exception
(addressed in \Cref{sec:reach}). \Cref{fig:mmcat} breaks the multimodal result down by task category, and
\Cref{fig:mlascale} reports the MLA decoupled-\kpe{} adapter fidelity and the transplant TTFT speedup
across models. \Cref{fig:adapters} diagrams the three attention-variant adapters of \Cref{sec:reach}: what
each representation caches, and exactly which slice of it the reposition touches.

\begin{figure}[tbp]
\centering
\begin{minipage}{0.46\linewidth}\centering
\includegraphics[width=\linewidth]{figA3_scorecard}
\caption{Composable agreement by model $\times$ content type.}
\label{fig:scorecard}
\end{minipage}\hfill
\begin{minipage}{0.52\linewidth}\centering
\includegraphics[width=\linewidth]{figA6_mla_scaling}
\caption{MLA adapter fidelity (a) and transplant TTFT speedup across models (b).}
\label{fig:mlascale}
\end{minipage}
\end{figure}

\begin{figure}[tbp]
\centering
\includegraphics[width=\linewidth]{figA5_multimodal}
\caption{Image-KV transplant by task category (perception / reasoning / agentic), full re-encode vs.\
transplant, across four vision-language models. The transplant tracks full accuracy category-by-category.}
\label{fig:mmcat}
\end{figure}

\begin{figure}[tbp]
\centering
\includegraphics[width=\linewidth]{figA7_adapters}
\caption{\textbf{The attention-variant adapters.} (a) MLA caches a position-free latent $c_t$ plus a small
decoupled-RoPE sub-vector \kpe{}; repositioning re-rotates only \kpe{} and copies the latent as-is. (b)
M-RoPE factors the rotary channels into temporal/height/width axes---sectioned (Qwen2.5-VL) or interleaved
(Qwen3-VL); moving an image within a trajectory re-rotates only the temporal channels, since the spatial
axes are intrinsic to the image. (c) Sliding-window layers default to a window-truncated cache, which
breaks any splice past $W$; keeping the full per-token KV and letting the attention mask enforce the
window restores uniform edit/splice semantics.}
\label{fig:adapters}
\end{figure}

\section{Long-horizon robustness and the cross-referential memory test}
\label{app:robustness}

This appendix collects two stress-tests for the leave-stale$+$erratum scheme. The first asks whether
per-edit errors \emph{compound} over a long agent trajectory (they do not); the second asks whether
splitting \emph{cross-referential} facts across independently-precompiled blocks breaks a two-hop chain---the
case \Cref{sec:memory}'s E4 left open (it does, and the mechanism explains why).

\paragraph{No compounding error over a long trajectory (\Cref{fig:horizon}).} A standing risk for any
leave-stale scheme is that small per-edit errors \emph{compound} over a long agent trajectory. We test it
directly. A single gated field (a clearance level) toggles between a granting and a denying value every turn
over a $28$-turn trajectory; we maintain ONE evolving KV cache---reuse the static prefix forever, apply each
state change as an appended erratum, never recompute the downstream---and compare its per-turn decision to a
\emph{full reprefill of the byte-identical token sequence} (errata included). The only difference between the
two is whether downstream KV was recomputed after each change, so per-turn agreement isolates exactly the
cost of leaving KV stale as a function of trajectory length. Across three families (Qwen3-8B, Llama-3.1-8B,
Mistral-7B) the next-token \emph{decision logits stay faithful}---cosine $0.987$--$0.999$ with a flat
first-third$\to$last-third profile (e.g.\ $0.992\to0.996$ on Llama)---and the agreement-vs-turn slope is
within $\pm0.01$/turn (no systematic decline). Discrete decision agreement is high in aggregate
($0.79$--$0.99$) but noisier than the cosine, because these gated decisions sit near the action boundary
(oracle accuracy $0.52$--$0.82$), where a sub-percent logit difference can flip a discrete choice---the same
boundary sensitivity noted for greedy CoT in \Cref{sec:memory}. The leave-stale$+$erratum cache does not
degrade with trajectory length: there is no compounding drift, only boundary noise.

\paragraph{Cross-referential facts (the \Cref{sec:memory} E4 test).} E4 found splitting \emph{independent}
relevant facts across independently-precompiled blocks decision-lossless. We test the case it left open---a
genuinely \emph{cross-referential} chain. A memory contains a \texttt{DEFINITION} line~$A$ naming which
setting governs the request, whose value~$B$ lives elsewhere; the decision needs the two-hop resolution
$A\!\to\!B$. We lay $A$ and $B$ so they straddle a block boundary and compare a transplant that \emph{splits}
them (boundary between $A$ and $B$) against one that keeps them \emph{colocated} (boundary moved past both),
holding everything else fixed. On Llama-3.1-8B---which performs this two-hop decision at full-recompute
accuracy $0.85$ ($n{=}80$, balanced)---splitting drops agreement with full recompute to $0.46$ versus $0.76$
colocated, a $0.30$ penalty ($25$ vs.\ $1$ discordant instances, McNemar $p{=}8\times10^{-7}$); an
\emph{independent} fact pair split the same way costs only $0.04$ (Mistral-7B). The mechanism explains it:
block~$B$ precompiled in isolation never attended to its referent~$A$, so the chain is not memoized and the
splice cannot restore it. (Qwen3-8B is at chance on the direct two-hop decision---consistent with
\Cref{sec:memory}'s observation that direct gated decisions are at chance for the reasoning-native family
$\le$32B---so it is uninformative here.) The guidance is simple: keep cross-referential facts within one
precompiled block; independent facts may be split freely.

\begin{figure}[tbp]
\centering
\includegraphics[width=\linewidth]{fig9_horizon}
\caption{\textbf{No compounding error over a long trajectory.} A gated field toggles every turn for $28$
turns; one evolving leave-stale$+$erratum cache vs.\ full reprefill of the identical text. (a) Per-turn
decision agreement stays high with boundary noise (no downward trend). (b) The decision-logit cosine stays
flat at $0.99{+}$---no drift with trajectory length.}
\label{fig:horizon}
\end{figure}

\begin{figure}[tbp]
\centering
\includegraphics[width=0.85\linewidth]{fig_e4_granularity}
\caption{\textbf{E4 --- granularity is a free knob} (\Cref{sec:memory}). Splitting memory into $S$
independently-precompiled blocks makes a localized edit $S\times$ cheaper and stays decision-lossless to
$S{=}16$; only genuinely cross-referential facts must be kept in one block (cross-referential test above).}
\label{fig:mem-e4}
\end{figure}

\section{The user-memory agent}
\label{app:memory-agent}
The agent of \Cref{sec:memory} keeps the layout \texttt{[system][trajectory][MEMORY][query]}.
The system prompt is prefilled once; the trajectory grows by cached deltas; the user-memory
chunk is precompiled once in isolation and RoPE-repositioned to float just before the query
each turn (an $O(L_\text{mem})$ re-rotation, no re-prefill) with one boundary token repaired;
a memory change is applied either by recompiling the isolated chunk ($O(L_\text{mem})$, once)
or by appending a salient erratum into the cached trajectory stream (append-only, composes
with prefix caching). Confirmatory protocol: gated decisions whose governing facts live in a
Markdown memory, balanced labels, chain-of-thought as the competent regime (direct one-shot
memory-gated decisions are at chance for $\le$8B models), with cluster-bootstrap CIs
($10^4$, persona-level), TOST equivalence ($\delta{=}0.03$; cosine $\ge0.98$), GEE-logistic
(cluster-robust), McNemar, and BH-FDR. Faithfulness is read against a token-matched
full-reprefill oracle. Pre-registered hypotheses, margins, and the inclusion gate (oracle
accuracy $\ge0.80$) are released with the code.

\section{Online serving and weight-editing methodology}
\label{app:systems}
\paragraph{Online vLLM serving (\Cref{fig:sys}a).} vLLM V1 \texttt{AsyncLLMEngine}, CUDA graphs enabled
(\emph{not} \texttt{enforce\_eager}), continuous batching, automatic prefix caching on, \texttt{bfloat16},
GPU memory utilization $0.85$. Workload: a shared ${\sim}8{,}066$-token policy (real
$\tau^2$-bench retail policy plus neutral padding) with one mutable field; $96$ requests per arm, $64$
output tokens each. Requests arrive as a \emph{Poisson} process at offered rates $\{2,4,8,16\}$ req/s and
unthrottled (saturation). We timestamp first-token and completion per request in the client loop (TTFT,
TPOT, end-to-end) and read each arm's APC hit-rate from the engine's Prometheus counters
(\texttt{vllm:gpu\_prefix\_cache\_hits}, \texttt{\ldots\_queries}). The in-prefix
baseline writes the new field value
early (invalidating downstream APC blocks); the erratum keeps the old prefix and appends the update.
Headline: throughput speedup grows $1.6\times{\to}14.5\times$ as load rises to saturation; $p90$ TTFT
$53$--$398\times$ lower; APC hit-rate $98.5\%$ vs.\ $1.0\%$.

\paragraph{Weight editing (\Cref{tab:weight}).} ROME is implemented from scratch (uncentered key covariance
$C{=}\mathbb{E}[kk^\top]$ at a mid MLP layer over a text sample, an optimized value $v^\star$, and the
closed-form rank-one update of \texttt{down\_proj}) and \emph{validated on the canonical factual edit}
(``the Eiffel Tower is in'' \emph{Paris}$\to$\emph{Rome}, locality preserved) before use, so the baseline
is faithful. LoRA fine-tunes ($r{=}8$, \texttt{q,v,down} projections) on the stale-context decision until it
flips. All methods are evaluated on the same gated decision; cross-request contamination is measured over
$8$ held-out orders that are genuinely still pending (correct $=$ cancel), and collateral over a $10$-item
battery of unrelated single-token gated decisions.

